% Lösungen Einheit 4 von 4 – Ungleichungen (zusammenbau v0.1)
\einheitenkopf{Einheit 4 von 4 · Ungleichungen}

\erg{20}{a) $x > 3$ \quad b) $x > 1$ (dort ist auch das Quadrat positiv) \quad c) $f(x) - g(x) = 0{,}5x^3 - 4x^2 + 8x = 0{,}5x \cdot (x - 4)^2$; Grenze $x = 0$, ausgenommen $x = 4$: $x > 0$ und $x \ne 4$ \quad d) $15k \cdot e^{-1{,}5} + 20 \ge 40 \Leftrightarrow k \ge \frac{20}{15 \cdot e^{-1{,}5}} \approx 5{,}98$; kleinster Wert $k \approx 5{,}98$ \quad e) $|p(x) - f(x)| \le 0{,}25$ aufstellen; $p(x) - f(x) = -0{,}25x^4 \le 0$, der Betrag ist $0{,}25x^4$; Randgleichung $0{,}25x^4 = 0{,}25$ lösen ($x = 1$); brauchbar für $0 \le x \le 1$.}
\erg{21}{a) $f(t) \ge 60$; Randstellen $t_1 \approx 4{,}95$, $t_2 \approx 12{,}10$; Zeitraum $[4{,}95; 12{,}10]$ min}
\erg{22}{a) Fehler: an der doppelten Nullstelle $x = 6$ ist das Produkt null, nicht positiv. Richtig: $x > 0$ und $x \ne 6$. \quad b) Ein Quadrat ist nie negativ; für $x \ne 7$ ist es positiv und lässt das Vorzeichen stehen, bei $x = 7$ ist das Produkt null. \quad c) Nullstellen $x_1 = 2$, $x_2 = 10$; $G(x) > 0$ für $2 < x < 10$, also zwischen 2000 und 10\,000 Stück \quad d) Randstellen $x_1 = -1$, $x_2 = 3$; $f(x) > g(x)$ für $-1 < x < 3$}
